\section*{ANEXO A – EXTRATO DA LEGISLAÇÃO BRASILEIRA SOBRE A GUARDA E PROTEÇÃO DE ANIMAIS DOMÉSTICOS}
\addcontentsline{toc}{section}{ANEXO A – EXTRATO DA LEGISLAÇÃO BRASILEIRA SOBRE A GUARDA E PROTEÇÃO DE ANIMAIS DOMÉSTICOS}

A seguir, apresentam-se os dispositivos legais federais citados ao longo deste trabalho que fundamentam a obrigação de restituição de animais encontrados e as penalidades aplicadas ao abandono e aos maus-tratos de cães e gatos no território brasileiro.

\vspace{0.5cm}
\noindent
\textbf{1. CÓDIGO CIVIL BRASILEIRO (LEI Nº 10.406, DE 10 DE JANEIRO DE 2002)}
\vspace{0.2cm}

\begin{quote}
\small
\textbf{Da Descoberta}

\textbf{Art. 1.233.} Quem quer que ache coisa alheia perdida há de restituí-la ao dono ou legítimo possuidor.
\par
\textit{Parágrafo único.} Não o conhecendo, o descobridor fará por encontrá-lo, e, se não o encontrar, entregará a coisa achada à autoridade competente.

\textbf{Art. 1.234.} Aquele que restituir a coisa achada, nos termos do artigo antecedente, terá direito a uma recompensa não inferior a cinco por cento do seu valor, e à indenização pelas despesas que houver feito com a conservação e transporte da coisa, se o dono não preferir abandoná-la.

\textbf{Art. 1.236.} A autoridade competente dará conhecimento da descoberta através da imprensa e outros meios de informação, somente expedindo editais se o seu valor os comportar.
\end{quote}

\vspace{0.5cm}
\noindent
\textbf{2. LEI FEDERAL Nº 14.064, DE 29 DE SETEMBRO DE 2020 (LEI SANSÃO)}
\vspace{0.2cm}

\begin{quote}
\small
Altera o art. 32 da Lei nº 9.605, de 12 de fevereiro de 1998, para aumentar as penas cominadas ao crime de maus-tratos aos animais quando se tratar de cão ou gato.

\textbf{Art. 1º} O art. 32 da Lei nº 9.605, de 12 de fevereiro de 1998, passa a vigorar acrescido do seguinte § 1º-A:

``\textbf{Art. 32.} [...]
\par
\textbf{§ 1º-A} Quando se tratar de cão ou gato, a pena para as condutas descritas no \textit{caput} deste artigo será de reclusão, de 2 (dois) a 5 (cinco) anos, multa e proibição da guarda.
\par
\textbf{§ 2º} A pena é aumentada de um sexto a um terço, se ocorre morte do animal.''
\end{quote}

\vspace{0.5cm}
\noindent
\textbf{3. SISTEMA NACIONAL DE CADASTRO DE ANIMAIS DOMÉSTICOS (SINPATINHAS)}
\vspace{0.2cm}

\begin{quote}
\small
Instituído no âmbito do Ministério do Meio Ambiente e Mudança do Clima (MMA) para integrar a gestão populacional de cães e gatos em território nacional, servindo de diretriz pública para a identificação individual e fomento à guarda responsável e aos programas municipais de controle de zoonoses e resgate.
\end{quote}

\clearpage
